\documentclass[aspectratio=169,11pt]{beamer}

% Match the visual language of the SIAM Waves talk.
\usetheme{Boadilla}
\usecolortheme{seahorse}
\useinnertheme{rectangles}
\useoutertheme{infolines}

\setbeamercolor{structure}{fg=blue!55!black}
\setbeamercolor{title}{fg=white,bg=blue!55!black}
\setbeamercolor{frametitle}{fg=blue!55!black}
\setbeamercolor{block title}{bg=blue!12,fg=blue!55!black}
\setbeamercolor{block body}{bg=blue!4}
\setbeamertemplate{navigation symbols}{}
\setbeamerfont{frametitle}{series=\bfseries,size=\large}
\setbeamerfont{title}{series=\bfseries,size=\Large}

\usepackage[T1]{fontenc}
\usepackage[utf8]{inputenc}
\DeclareFontShape{T1}{cmss}{b}{n}{<->ssub*cmss/bx/n}{}
\usepackage{amsmath, amssymb, bm}
\usepackage{graphicx}
\usepackage{xcolor}
\usepackage{hyperref}
\hypersetup{colorlinks=true,urlcolor=blue!55!black}

% ---------- Macros ----------
\newcommand{\R}{\mathbb{R}}
\newcommand{\uu}{\bm{u}}
\newcommand{\pp}{p}
\newcommand{\Rey}{\mathrm{Re}}
\newcommand{\nuu}{\nu}
\newcommand{\tmark}[1]{\hfill{\scriptsize\textcolor{gray}{#1}}}

% Uncomment if you want speaker notes:
% \usepackage{pgfpages}
% \setbeameroption{show notes on second screen=right}

\title[Turbulence: canonical flows]{Turbulence through canonical flows}
\subtitle{Transition, spectra, examples, data, and computation}
\author[Miller]{Laura A. Miller}
\institute[UArizona]{Department of Mathematics, University of Arizona}
\date{}

\begin{document}
% -------------------------
\begin{frame}[plain]
  \begin{beamercolorbox}[sep=10pt,center]{title}
    {\bfseries\Large Turbulence through canonical flows}\\[3pt]
    {\normalsize Transition, spectra, examples, data, and computation}
  \end{beamercolorbox}
  \vspace{1.35em}
  \begin{columns}[T,onlytextwidth]
  \begin{column}{0.55\textwidth}
    \textbf{\large Laura A. Miller}\\[4pt]
    Department of Mathematics\\
    University of Arizona\\[12pt]
    {\small IBFE Tutorials}
  \end{column}
  \begin{column}{0.45\textwidth}
    \begin{block}{Guiding idea}
    Canonical flows isolate mechanisms that recur in more complicated systems.
    \end{block}
  \end{column}
  \end{columns}
  \vfill
  {\scriptsize Adapted from a \href{https://cassyni.com/events/HsafxyMiJvA1JVCU9eef3T/seminar}{lecture by Steven L. Brunton}.}
  \\[2pt]{\tiny AI tools were used to assist with formatting, organization, and clarification of the instructor's original notes.}
\end{frame}

% -------------------------
\begin{frame}{Lecture roadmap}
\begin{itemize}
  \item Why turbulence is hard (and why it matters)
  \item Canonical flows: a ``vocabulary'' for turbulence
  \item Examples: pipe, wake, boundary layer, cavity, jets, mixing layers
  \item ``More complex'' flows (geometry, separation, multi-scale physics)
  \item Open-source resources: databases and CFD codes
\end{itemize}
\end{frame}


% -------------------------
\begin{frame}{Why turbulence is hard}
\begin{columns}[T,onlytextwidth]
\column{0.56\textwidth}
\begin{itemize}
  \item Turbulence is a \emph{multiscale}, nonlinear dynamics problem whose observed state can depend on disturbance history.
  \item Small perturbations can grow, interact, and reorganize the flow.
  \item Many flows are \emph{linearly stable} yet become turbulent through finite-amplitude disturbances.
  \item Predicting transition and statistics (drag, mixing, heat transfer, intermittency) is still challenging.
\end{itemize}

\column{0.44\textwidth}
\includegraphics[width=\textwidth]{pictures/5-1/turbulence.png}

\vspace{0.35em}
{\scriptsize Source: to be confirmed.}

{\footnotesize The \emph{Feynman Lectures} (Vol.\ II, Ch.~41) give a classic discussion of why real viscous flows are difficult to predict.}
\end{columns}
\end{frame}

% -------------------------
\begin{frame}{Canonical flows (Van Dyke ``album'')}
\begin{columns}[T,onlytextwidth]
\column{0.55\textwidth}
\begin{block}{Idea}
A small set of ``canonical'' configurations show up \emph{everywhere}.  Learn these and you gain a language for:
\begin{itemize}
  \item transition mechanisms (instability \(\rightarrow\) breakdown),
  \item coherent structures (vortices, streaks, shear layers),
  \item scaling laws and nondimensional groups.
\end{itemize}
\end{block}
\column{0.4\textwidth}
\includegraphics[width=0.75\textwidth]{pictures/5-1/album.png}

{\scriptsize M. Van Dyke, \emph{An Album of Fluid Motion} (Parabolic Press, 1982), cover.}
\end{columns}
\end{frame}

% -------------------------
\begin{frame}{A minimal tool: Reynolds number}
\begin{itemize}
  \item For chosen velocity and length scales, Reynolds number compares inertia to viscosity:
  \[
    \mathrm{Re}=\frac{UL}{\nu}.
  \]
  \item ``High'' \(\mathrm{Re}\) often permits thinner shear layers and wider scale separation, but does not guarantee turbulence.
  \item Treat \(\mathrm{Re}\) as a control parameter, not a universal one-to-one label for complexity:
  geometry, forcing, roughness, and disturbance history matter.
\end{itemize}
\end{frame}

% =========================
% Canonical flows
% =========================

\section{Canonical flows}

% -------------------------
\begin{frame}{Pipe flow}
\begin{columns}[T,onlytextwidth]
\column{0.5\textwidth}
\begin{itemize}
  \item Classic internal flow: long straight pipe, pressure-driven.
  \item The Hagen--Poiseuille solution exists for every \(\mathrm{Re}\) and is linearly stable (proved for axisymmetric disturbances; computed stable to \(\mathrm{Re}=10^{7}\), Meseguer \& Trefethen 2003), yet finite-amplitude disturbances can trigger turbulence.
  \item Transition often involves localized ``puffs'' and ``slugs'' (intermittency).
  \item Active control is possible: carefully designed perturbations can suppress turbulence (drag reduction).
\end{itemize}
\column{0.45\textwidth}
\includegraphics[width=\textwidth]{pictures/5-1/t_0300.jpg}

{\scriptsize Source: to be confirmed.}
\end{columns}
\end{frame}

%-------------------------------
\begin{frame}{Pipe flow: Slugs and puffs}
\begin{itemize}
\item Transitional pipe flow can coexist with laminar regions and is intrinsically intermittent.
\item Localized turbulent structures have distinct dynamics:
\begin{itemize}
\item \textbf{puffs} remain localized while advecting and may decay or split stochastically;
\item at higher \(\mathrm{Re}\), \textbf{slugs} expand by turbulent invasion at their fronts.
\end{itemize}
\item The crossover is a range of behavior rather than a single universal threshold.
\end{itemize}

\centering
\includegraphics[width=0.7\textwidth]{pictures/5-1/turbulence-pipe.png}

\vspace{0.35em}
{\footnotesize Source: Graham, M. Turbulence spreads like wildfire. Nature 526, 508–509 (2015).}
\end{frame}

%-------------------
\begin{frame}{Puffs, slugs, and spectral similarity}
\begin{columns}[T,onlytextwidth]
\column{0.54\textwidth}
\begin{itemize}
\item The image sketches spectra sampled from the marked pipe-flow regions.
\item For high-$\Rey$, locally isotropic \textbf{3D} turbulence, Kolmogorov similarity hypothesizes an inertial range
\[
E(k)=C_K\varepsilon^{2/3}k^{-5/3}.
\]
\item This is asymptotic and statistical---not exact universality. Intermittency, anisotropy, walls, forcing, and finite $\Rey$ produce deviations.
\item A 2D inverse-energy cascade may also have a $-5/3$ range, but its transfer direction and mechanism differ; enstrophy transfers downscale.
\end{itemize}
\column{0.43\textwidth}
\centering
\includegraphics[width=\textwidth]{pictures/5-1/1-goingwiththe.jpg}

\vspace{0.25em}
{\scriptsize Source image: \href{https://phys.org/news/2020-01-insights-mysterious-fluid-motions.html}{Phys.org (2020)}.}
\end{columns}

\end{frame}

% -------------------------
\begin{frame}{Wake flows}
\begin{columns}[T,onlytextwidth]
\column{0.5\textwidth}
\begin{itemize}
  \item Flow past a body leaves a \emph{wake} with a velocity deficit and strong shear.
  \item Bluff bodies can shed a K\'arm\'an vortex street, producing unsteady loading and noise; lifting surfaces also generate trailing and tip vortices.
  \item Wakes are central to:
  \begin{itemize}
    \item aircraft wing tips and propellers,
    \item wind farms and arrays,
    \item animal locomotion and swimming.
  \end{itemize}
\end{itemize}
\column{0.45\textwidth}
\includegraphics[width=\textwidth]{pictures/5-1/t_0360.jpg}

{\scriptsize Source: to be confirmed.}
\end{columns}
\end{frame}


% -------------------------
\begin{frame}{Boundary layer}
\begin{columns}[T,onlytextwidth]
\column{0.56\textwidth}
\begin{itemize}
  \item Near a wall, no-slip forces a thin region of strong shear: the boundary layer.
  \item Key questions:
  \begin{itemize}
    \item when does laminar flow transition to turbulence?
    \item how do near-wall streaks and vortices sustain turbulence?
    \item how do roughness and pressure gradients change drag?
  \end{itemize}
  \item Boundary layers dominate drag and heat transfer on vehicles, buildings, and animals.
\end{itemize}
\column{0.44\textwidth}
\includegraphics[width=\textwidth]{pictures/5-1/turbulent-boundary.png}

{\scriptsize Smoke visualization of a turbulent boundary layer; Van Dyke, \emph{An Album of Fluid Motion} (turbulent boundary layer plate; also reproduced in Dou 2018).}
\end{columns}
\end{frame}

% -------------------------
\begin{frame}{Boundary layer: spatial development}
\centering
\includegraphics[width=0.92\textwidth]{pictures/5-1/boundary-layer.png}
\vspace{0.35em}

{\scriptsize Boundary-layer schematic after GRAS Sound \& Vibration, \emph{The boundary layer---laminar and turbulent flows} (blog; source to confirm), \url{https://www.grasacoustics.com/blog/aerodynamic-measurements-components-of-the-boundary-layer}.}
\end{frame}

% -------------------------
\begin{frame}{Cavity flow}
\begin{columns}[T,onlytextwidth]
\column{0.58\textwidth}
\begin{itemize}
  \item In an \textbf{open engineering cavity}, a shear layer spans the opening; pressure/acoustic feedback can create resonant oscillations.
  \item An intracranial aneurysm is instead a \textbf{cavity-like recirculating flow}: pulsatility, geometry, and vessel-wall conditions shape its hemodynamics.
  \item Related applications:
  \begin{itemize}
    \item aneurysm inflow jets and recirculation,
    \item gaps, steps, and surface defects,
    \item flow-induced vibration.
  \end{itemize}
\end{itemize}
\column{0.42\textwidth}
\includegraphics[width=\textwidth]{pictures/5-1/aneurysm.png}
{\footnotesize Source: Boniforti MA, Magini R, Orosco Salinas T. Hemodynamic Investigation of the Flow Diverter Treatment of Intracranial Aneurysm. Fluids. 2023; 8(7):189. }
\end{columns}
\end{frame}



% =========================
% More complex flows and resources
% =========================

% -------------------------
\begin{frame}{Jets, mixing layers, and separated flows (overview)}
\begin{itemize}\itemsep4pt
  \item \textbf{Round jet:} downstream of the potential core the mean profile becomes \emph{self-similar}, the half-width grows roughly linearly with a spreading angle of order $10$--$12^\circ$, and Kelvin--Helmholtz roll-up of the annular shear layer seeds the initial vortex rings.
  \item \textbf{Plane mixing layer:} two streams of different speed produce the Brown--Roshko coherent spanwise rollers, which pair and merge; the growth rate scales with the velocity ratio $(U_1-U_2)/(U_1+U_2)$ rather than with $\mathrm{Re}$.
  \item \textbf{Separated flows:} a backward-facing step separates at a fixed geometric edge; at high $\mathrm{Re}$ the mean reattachment length is about $6$--$7$ step heights, and the shear layer, recirculation bubble, and reattachment zone each have their own statistics.
  \item \textbf{Combinations:} real configurations superpose these---a jet issuing into a cavity, a wake feeding a mixing layer, a step downstream of a boundary layer.
  \item \textbf{Multi-scale physics:} geometry sets the large scales; viscosity sets the small ones. The gap between them is what makes prediction expensive.
\end{itemize}
\end{frame}

% -------------------------
\begin{frame}{Open data: turbulence databases}
\begin{columns}[T,onlytextwidth]
\column{0.57\textwidth}
\begin{itemize}
  \item Public datasets let you test ideas without running a massive DNS yourself.
  \item \href{https://turbulence.pha.jhu.edu/}{Example: Johns Hopkins Turbulence Databases (JHTDB)}
  \begin{itemize}
    \item isotropic turbulence, channel flow, boundary layers, geophysical, \ldots
    \item query-able via web/API-style interfaces
    \item great for teaching statistics, structure functions, spectra, and coherent-structure detection
  \end{itemize}
\end{itemize}
\column{0.43\textwidth}
\includegraphics[width=\textwidth]{pictures/5-1/JHU-dataset.png}

{\scriptsize Johns Hopkins Turbulence Databases, \href{https://turbulence.pha.jhu.edu/}{turbulence.pha.jhu.edu}.}
\end{columns}
\end{frame}

% -------------------------
\begin{frame}{Open-source CFD codes}
\begin{columns}[T,onlytextwidth]
\column{0.52\textwidth}
\begin{block}{Two common ecosystems}
\begin{itemize}
  \item \href{https://openfoam.org/}{\textbf{OpenFOAM}}: open-source finite-volume framework with many RANS/LES solvers.
  \item \href{https://nek5000.mcs.anl.gov/}{\textbf{Nek5000}}: high-order spectral-element code; widely used for DNS/LES in complex geometries.
\end{itemize}
\end{block}
\begin{itemize}
  \item Choosing a tool depends on: accuracy needs, geometry, turbulence model, and HPC resources.
\end{itemize}
\column{0.44\textwidth}
\includegraphics[width=\textwidth]{pictures/5-1/nek5000-1000x500.jpg}

{\scriptsize Nek5000 project banner image, Argonne National Laboratory, \href{https://nek5000.mcs.anl.gov/}{nek5000.mcs.anl.gov}.}
\end{columns}
\end{frame}

% -------------------------
\begin{frame}{Takeaways}
\begin{itemize}
  \item Turbulence is everywhere, but canonical flows give structure to the chaos.
  \item Each canonical flow highlights different mechanisms (instability, shear, walls, separation, resonance).
  \item Open videos, datasets, and codes make it easier than ever to \emph{see} turbulence and to explore it quantitatively.
\end{itemize}

\vspace{0.8em}

\end{frame}

\begin{frame}{Selected references and source notes}
\scriptsize
\begin{itemize}
  \item R. P. Feynman, R. B. Leighton, and M. Sands, \emph{The Feynman Lectures on Physics}, Vol.~II, Ch.~41, \href{https://www.feynmanlectures.caltech.edu/II_41.html}{``The Flow of Wet Water''}.
  \item K. Avila, D. Moxey, A. de Lozar, M. Avila, D. Barkley, and B. Hof, ``The onset of turbulence in pipe flow,'' \emph{Science} \textbf{333} (2011), 192--196. \href{https://doi.org/10.1126/science.1203223}{doi:10.1126/science.1203223}.
  \item A. N. Kolmogorov, ``The local structure of turbulence in incompressible viscous fluid for very large Reynolds numbers,'' \emph{Proc. R. Soc. A} \textbf{434} (1991 translation), 9--13. \href{https://doi.org/10.1098/rspa.1991.0075}{doi:10.1098/rspa.1991.0075}.
  \item R. H. Kraichnan, ``Inertial ranges in two-dimensional turbulence,'' \emph{Phys. Fluids} \textbf{10} (1967), 1417--1423. \href{https://doi.org/10.1063/1.1762301}{doi:10.1063/1.1762301}.
  \item H. Schlichting and K. Gersten, \emph{Boundary-Layer Theory}, 9th ed., Springer (2017).
  \item Official resources: \href{https://turbulence.pha.jhu.edu/}{JHTDB}, \href{https://openfoam.org/}{OpenFOAM Foundation}, and \href{https://nek5000.mcs.anl.gov/}{Nek5000}.
\end{itemize}
\end{frame}
\end{document}
